\subsection{Nouvelles disciplines d'étude}

%% Introduction au problèmes %%

Pour finir, les profils des chercheurs ont changé au contact de la numérisation, ils n'ont pas les mêmes besoins qu'avant qu'elle n'intègre les salles de lectures. Le premier indicateur en est l'intérêt que portent les chercheurs aux archives audiovisuelles, qui comme nous l'avons vu, n'a pas toujours bénéficié de la même légitimité que l'étude des manuscrits ou des archives. Ce constat s'applique également à une institution patrimoniale plus ancienne comme la BNF, qui au cours d'une études sur les usages documentaire des chercheurs détermine que l'intérêt de leurs lecteurs se sont de plus en plus portés vers de nouvelles disciplines\footcite{pardé2015}. Si l'on considère que l'indexation est établit en fonction des usages et besoins de ses utilisateurs au moment où elles sont utilisées, cela crée aujourd'hui un décalage entre l'indexation utilisée dans les notices documentaires et ce que recherchent effectivement les usagers. 

%% Application à chaque public %%

Il devient alors nécessaire de comprendre quelle indexation a été utilisée et de développer les pratiques nécessaires pour contourner ces problématiques. Sans cela, une recherche par mot-clé n’aboutis pas au résultat espéré, ils sont amoindri car le vocabulaire utilisé est obsolète, pensons par exemple au termes de \enquote{féminicide}, \enquote{réchauffement climatique} qui ne trouveront pas leur place dans des notices anciennes. Impossible dans ce cas, en faisant uniquement appel l'indexation d'analyser le traitement de l'actualité sur ces sujets dans les journaux télévisés des années 80. La clef de ce décalage est bien souvent l'expertise documentalistes, qui dotés de leurs expertises sur les collections ont connaissance de ce décalage et peuvent offrir des solutions ou même éventuellement mettre en place une médiation et construire des corpus en fonction de l'actualité comme c'est réalisé sur le site ina.fr mais également dans les émissions produites par l'INA, étudiant des sujets d'actualités à la lumières de documents audiovisuels d'archives. 

\subsection{De nouveaux publics}

Les publics touchés par ces collections sont également bouleversés, le grand public est de plus en plus touché par ces collections mais il ne les consomme pas en salle de lecture comme le font les chercheurs. La consommation des collections s'est déplacée en ligne. En conséquent les habitudes de consommation s'adaptent aux formats conçus à ce grand public des sites web où la découverte fortuite de nouvelles images est aisée à l'époque des algorithmes personnalisés. Au sein de la salle de lecture, cette aspect de découverte se perd. Difficile d'arriver en salle de lecture sans savoir ce que l'on recherche, les outils de consultation ne sont conçus que pour répondre à une chaîne de caractères précise, pas à une vague envie de découvrir son patrimoine audiovisuel. Là encore, il n'y a que les documentalistes pour venir en aide à ces nouveaux publics s'aventurant en salle de lecture. 